\scanpage{95}

\reversemarginpar
\margin{\begin{tikzpicture}[x=1cm,y=1cm,baseline=(current bounding box.north)]
\draw[line width=.25pt] (.04,0) .. controls (.00,-.20) and (.03,-1.1) .. (.10,-2.4) .. controls (.17,-3.6) and (.31,-5.2) .. (.28,-5.73);
\draw[line width=.25pt] (.28,-5.79) .. controls (.10,-5.81) and (-.05,-5.96) .. (-.24,-6.13);
\end{tikzpicture}}
\normalmarginpar

que $N(x)$ induit sur $R^n$ la fonction $N_0(x)$, et qu'on a $N(\lambda x)=|\lambda|N(x)$ et $N(x)=0\Rightarrow x=0$, il faut seulement prouver $N(x+y)\leq N(x)+N(y)$ (\textcolor{red}{\uline{\textcolor{black}{est}}} valable $=$ pour $x,y\in c_0$.) Pour ceci, on \deleted{\unsure{regarde} $(x_i)$ et $(y_i)$ comme \unsure{les}}\inserted{remarque que si} $u_x$ est l'opération de multiplication par $x$ dans $\ell^2$ (qui est un opérateur cpt pour $x\in c_0$), on a \deleted{$u_{x+y}=u_x$} $\rho_i(u_x)=\rho_i(x)$, d'autre part \deleted{\unsure{que|$\rho_i$}} $u_{x+y}=u_x+u_y$, et l'inégalité à prouver résulte alors du N\textsuperscript{\uline{o}}~4, th.~4, corollaire. En $=$ temps, les normes \deleted{\unsure{ft}}\inserted{de Schatten} ainsi associées à des normes \deleted{\unsure{de}} de Schatten $N_0$ sur $R^n$ \deleted{donnent} démontrent qu'il existe une infinité \inserted{continue} de normes de Schatten sur $R^{(\underline N)}$ équivalentes à la norme induite par $c_0$.

Soit $N$ \deleted{\unsure{est}} une norme de Schatten sur \deleted{$\unsure{R^I|R^n}$ \unsure{finies|fini}} $R^{(\underline N)}$. \deleted{Soit $\unsure{\ell_N}$ $\ell^N$ \unsure{l'espace des}}\inserted{Pour \deleted{toute}} suite \deleted{\unsure{réelle}} complexe $x=(x_i)$ \deleted{telle que \unsure{la suit}} \deleted{\unsure{on pose|$N\quad\rho$}} \deleted{soit $\unsure{N(x\mathord{-}x|N(x\,x}$}

\[
\deleted{N(|x_1|,\mathord{\cdot}\,\mathord{-}|x_n|,0,0,\mathord{-}\,\mathord{-}\,\mathord{\cdot})\quad\text{croît avec }n,\quad\text{on pose.}}
\]

Soit $|x|_n=(|x_1|,\mathord{\cdot}\,\mathord{-}|x_n|,0,\mathord{-}\,\mathord{-}0)$. \deleted{\unsure{On a aussitôt|Comme on a}}\inserted{Si $n$ varie} \deleted{alors $|x+y|_n\leq |x|_n+|y|_n$} $|x|_n$ est une suite croissante d'éléments positifs de $R^{(N)}$, posons

\[
\text{(4)}\qquad N(x)=\lim_n N(|x|_n)
\]

\inserted{c'est un nb $\geq 0$ fini ou infini} La valeur \deleted{\unsure{impropre|infinie|ainsi} \unsure{est définie|désignée}} de $N(x)$, qui coïncide avec \deleted{\unsure{$N(x)$|celle}} \inserted{quand $x\in R^{(N)}$.} On a

\[
\text{(5)}\qquad N(\lambda x)=|\lambda|N(x),\qquad N(u+y)\leq N(\unsure{u|x})-N(y)
\]

\deleted{si $N(x)<+\infty$} \inserted{et dans} la première formule, \deleted{\unsure{il est à dire|étant dit} que \unsure{son deuxième mb a un sens}, i.e $|\lambda|$} \inserted{il faut prendre au deuxième}
\reversemarginpar
\margin{\textcolor{red}{?}}
\normalmarginpar
mb la valeur \textcolor{red}{\uline{\textcolor{black}{zéro}}} s'il \deleted{\unsure{est sous}}\inserted{se présente sous la} forme indéterminée $0\times\infty$. La deuxième formule résulte de ce que \deleted{$N(x$} $|x+y|_n\leq |x|_n+|y|_n$, d'où $N(|x+y|_n)\leq N(|x|_n+|y|_n)\leq N(|x|_n)+N(|y|_n)$.

\deleted{Soit} Soit $\ell^N$ l'espace des suites complexes $x$ telles que $N(x)<\infty$, espace \unsure{contenant|vectoriel} au moins \deleted{\unsure{$C^{(N)}$|celles}} les suites complexes nulles à partir \deleted{\unsure{de rang|de mult}} d'un certain rang. \unsure{Sur|Par|Pour} $(5)$, on voit que $N(x)$ est une \uline{norme} sur $\ell^N$.

\deleted{\unsure{De plus|D'ailleurs}, $\ell^N$ est \unsure{un espace de} \illegible{2} \unsure{positivité}. D'où \unsure{les inj} $R^{(N)}\subset\ell^N$, et $N(x)$ induit \unsure{par hypothèse}, \unsure{donc} ($=$ \illegible{1} $\unsure{\alpha^i|\alpha^{(i)}}$), est \illegible{2} \unsure{le candidat|la convergence} dans $\ell^N$, \unsure{elles coïncident donc forcément|la convergence en chaque point}}

\reversemarginpar
\margin{96}
\normalmarginpar
